\section{Architecture \& Environment}
\begin{frame}{Architecture \& Environment}
  \begin{itemize}
    \item Project Architecture
    \item Framework and Libraries
    \item Hardware
  \end{itemize}
\end{frame}

\begin{frame}{Project Architecture \textendash\ Source Files}
  The project is organized in different \textbf{Python} scripts in the root folder, these scripts can either be \textit{entrypoints} or \textit{utility} scripts

  \begin{itemize}
    \item \texttt{baseline.py}: Entrypoint model training for~\ref{sec:task1}.
    \item \texttt{degraded.py}: Entrypoint model training for~\ref{sec:task2}.
    \item \texttt{oneshot.py}: Entrypoint model training for~\ref{sec:task3}.
    \item \texttt{evaluate.py}: Entrypoint inferenza per tutti i modelli.
    \item \texttt{models.py}: Utility file containing all used models.
    \item \texttt{transforms.py}: Utility file containing some transforms used for image enhancement / augmentation.
    \item \texttt{utils.py}: Utility containing other functions used globally in the project.
  \end{itemize}
\end{frame}

\begin{frame}{Project Architecture \textendash\ Other Folders}
  \begin{itemize}
    \item \texttt{data}: Contains the dataset subdivided in \texttt{test, test\_degradato, train, valid} and in 100 different folders (1 per class)
    \item \texttt{docs}: Contains documentation and references about the project code and dataset.
    \item \texttt{models}: Contains \textbf{pytorch model checkpoints}
    \item \texttt{report}: Contains this report presentation.
  \end{itemize}
\end{frame}

\begin{frame}{Project Libraries \& Frameworks}
  \begin{itemize}
    \item \texttt{matplotlib}: MATLAB-like plotting library
    \item \texttt{numpy}: MATLAB-like matrices and numerical operations
    \item \texttt{torch}: PyTorch main Machine Learning framework used
    \item \texttt{torchvision}: Pytorch sub-library for computer vision
    \item \texttt{torchinfo}: model summary and number of parameters
    \item \texttt{timm}: pre-trained and famous model architectures sourcing
    \item \texttt{scikit-learn}: some utility functions to compute metrics on inference
    \item \texttt{opencv}: open computer vision library used in task3 to apply handcrafted feature extraciton and other operations
  \end{itemize}
\end{frame}

\begin{frame}{ROCm and GPUs}
  In order to train the models we used \textbf{ROCm} (AMD's open source high performace compute library) on 2 AMD consumer cards: RX6950 \& RX9070.
  Significant effort to setup a working environment.

  Final project env setup allows both CPU and GPU-accellerated conditional setups, by conditionally sourcing libraries from 3 different sources:
  \begin{itemize}
    \item \textbf{CPU}:\@ \texttt{uv sync --extra cpu}
    \item \textbf{GPU RX6950}:\@ \texttt{uv sync --extra stable}
    \item \textbf{GPU RX9070}:\@ \texttt{uv sync --extra nightly}
  \end{itemize}
\end{frame}